\section*{Chapter \ref{ch:g10:ricardo-model} --- Ricardo's Model of Comparative Advantage}
% figdata: figdata/grade-10/03-ricardo-model/out/ (wages_selected, wages_summary, wages_productivity);
% Brenland and Kestria hours: code/shared/country/ricardo.csv (invented).

\begin{solution}{exo:g10:ricardo-model:1}
Brenland has the absolute advantage in both goods (4 hours against 6 for cloth, 2 against 12 for grain). Absolute
advantage compares the countries good by good; trade depends on what each hour can earn in one good compared with
the other inside each country, that is on the ratios $a_C/a_G$, the comparative advantage. A country better at
everything still gains by moving its hours to the good where its lead is largest.
\end{solution}

\begin{solution}{exo:g10:ricardo-model:2}
Brenland: $2/4 = 0.5$ roll of cloth per tonne of grain; Kestria: $12/6 = 2$ rolls per tonne. The \omterm{def:g10:ricardo-model:autarky-price}{autarky prices} of
cloth were 2 and 0.5 tonnes per roll: $1/2 = 0.5$ and $1/0.5 = 2$, the inverses.
\end{solution}

\begin{solution}{exo:g10:ricardo-model:3}
In A a bicycle costs $20/5 = 4$ chairs; in B, $30/15 = 2$ chairs. B has the comparative advantage in bicycles (and
A in chairs), although A has the absolute advantage in both. The world price must lie between 2 and 4 chairs per
bicycle.
\end{solution}

\begin{solution}{exo:g10:ricardo-model:4}
Reading the cloud: at \usd{10000} of output per worker, compensation of roughly \usd{4000} to \usd{5000}; at
\usd{100000}, roughly \usd{50000} to \usd{60000}. (The line through the cloud gives about \usd{4500} and
\usd{50800}.)
% figdata: wages_productivity (fitted line computed from the table: 4483 and 50779)
\end{solution}

\begin{solution}{exo:g10:ricardo-model:5}
What matters is not hours but money costs, and wages differ: Kestria's wage is lower than Brenland's because its
workers produce less. With a wage between 1.5 and 6 times lower, Kestria's cloth, where it is only 1.5 times less
productive, costs less in money than Brenland's. So Kestria sells cloth to Brenland, and both gain.
\end{solution}

\begin{solution}{exo:g10:ricardo-model:6}
Brenland gets a roll for 1.5 tonnes of grain, $1.5 \times 2 = 3$ hours, instead of 4: it saves 1 hour per roll
(2 at $p = 1$). Kestria gets a tonne for $1/1.5$ roll, $6/1.5 = 4$ hours, instead of 12: it saves 8 hours per tonne
(6 at $p = 1$). Kestria's gain rose because the price of its export, cloth, rose: its \omterm{def:g10:ricardo-model:terms-of-trade}{terms of trade} improved, and
Brenland's worsened.
\end{solution}

\begin{solution}{exo:g10:ricardo-model:7}
With 3 hours, Kestria's \omterm{def:g10:ricardo-model:autarky-price}{autarky price} of cloth becomes $6/3 = 2$ tonnes, the same as Brenland's: neither country has
a comparative advantage, and trade brings no gain (the only price at which both would trade is 2, where trading
saves no hours). With 2 hours, Kestria's \omterm{def:g10:ricardo-model:autarky-price}{autarky price} of cloth becomes $6/2 = 3$, above Brenland's 2: Kestria now
has the comparative advantage in grain and Brenland in cloth. The pattern reverses, at a world price between 2 and 3.
\end{solution}

\begin{solution}{exo:g10:ricardo-model:8}
$w_B/w_K = 3/1.5 = 2$. In units of Kestria's hourly wage: a roll of cloth costs 6 in Kestria and $4 \times 2 = 8$ in
Brenland, so Kestria sells cloth more cheaply; a tonne of grain costs 12 in Kestria and $2 \times 2 = 4$ in Brenland, so
Brenland sells grain more cheaply.
\end{solution}

\begin{solution}{exo:g10:ricardo-model:9}
Kestria makes 200 rolls and sells each for 1.5 tonnes: it can consume any $(C, G)$ with $G = 1.5\,(200 - C)$, that is
$1.5\,C + G = 300$. Keeping 100 rolls, it gets $1.5 \times 100 = 150$ tonnes of grain. Without trade, 100 rolls take 600
hours, and the other 600 make $600/12 = 50$ tonnes.
\end{solution}

\begin{solution}{exo:g10:ricardo-model:10}
The ratio of output per worker is $139680/19440 \approx 7.19$, and the ratio of compensation
$82660/10370 \approx 7.97$. Labour's share: $82660/139680 \approx
0.59$ in the United States and $10370/19440 \approx 0.53$ in India. Pay differs by a little more than output, because
the share paid to labour is a little higher in the United States.
\end{solution}

\begin{solution}{exo:g10:ricardo-model:11}
\omterm{def:g10:ricardo-model:autarky-price}{Autarky}: Brenland makes $600/4 = 150$ rolls and $600/2 = 300$ tonnes; Kestria $600/6 = 100$ rolls and $600/12 = 50$
tonnes. World: 250 rolls and 350 tonnes. With trade: Kestria makes $1200/6 = 200$ rolls; Brenland makes the other 50
rolls with 200 hours and grain with its other \num{1000} hours: 500 tonnes. World: 250 rolls and 500 tonnes. World
grain rises by 150 tonnes with the same cloth.
\end{solution}

\begin{solution}{exo:g10:ricardo-model:12}
The model is contradicted if a country exports the good in which its relative labour requirement is \emph{higher}:
for example, Kestria exporting grain to Brenland while needing 6 times Brenland's hours for grain and only 1.5 times
for cloth. Reasons outside the model: other factors of production (Kestria may have much good land, which the model
ignores), transport costs or tariffs that differ between goods, differences of quality inside a ``good'', or hours
measured badly.
\end{solution}

\begin{solution}{exo:g10:ricardo-model:13}
Kestria's hours over Brenland's: tools $5/5 = 1$, cloth $6/4 = 1.5$, grain $12/2 = 6$. Kestria makes a good more
cheaply when $a^K w_K < a^B w_B$, that is when $a^K/a^B < w_B/w_K$. If Kestria made grain more cheaply, $w_B/w_K > 6$,
which is larger than 1.5 and 1, so Kestria would make all three goods more cheaply and Brenland could sell nothing.
Brenland could then not pay for any import, its workers would be idle and its wage would fall: trade must be paid for
with exports, so this cannot last. If Brenland makes cloth more cheaply, $w_B/w_K < 1.5$, which is less than 6, so
Brenland also makes grain more cheaply. With $1 < w_B/w_K < 1.5$: Brenland makes cloth and grain, Kestria makes tools
(its hours ratio, 1, is below the wage ratio). The goods line up as a chain, and the wage ratio decides where the chain
is cut.
\end{solution}

\begin{solution}{exo:g10:ricardo-model:14}
At $p = 1$, $w_B = 3\,w_K$. A tonne of grain costs $2 \times 3 = 6$ Kestria wage-hours in Brenland against 12 in
Kestria: Brenland's lead in grain (six times the output per hour) is larger than its wage gap (three times), so it
still sells grain more cheaply. With a wage seven times Kestria's, a tonne of Brenland grain would cost $2 \times 7 =
14 > 12$, and a roll of cloth $4 \times 7 = 28 > 6$: Brenland could export nothing. Its workers would have no buyers
abroad, and its wage would have to fall back below six times Kestria's.
\end{solution}

\begin{solution}{exo:g10:ricardo-model:15}
Let $p_B$ and $p_K$ be the \omterm{def:g10:ricardo-model:autarky-price}{autarky prices} of cloth, with $p_K < p_B$. If $p > p_B$, then $p$ is above both, and by
\cref{prop:g10:ricardo-model:pattern} both countries make only cloth: nobody makes grain, so there is no grain to import
and no trade in which each exports one good. If $p < p_K$, both make only grain, for the same reason. So each country
can export a different good only if $p_K \le p \le p_B$; strictly inside, both specialise completely, and at an end one
country may make both goods.
\end{solution}

\begin{solution}{pb:g10:ricardo-model:1}
Quantities are in millions of rolls and millions of tonnes.

\textbf{1.} Brenland: $4C + 2G = 600$; at most 150 rolls (no grain) or 300 tonnes (no cloth). Kestria: $6C + 12G =
1800$; at most 300 rolls or 150 tonnes.

\textbf{2.} Brenland: $4/2 = 2$ tonnes per roll. Kestria: $6/12 = 0.5$ tonne per roll.

\textbf{3.} Brenland: 300 hours on each good, $300/4 = 75$ rolls and $300/2 = 150$ tonnes. Kestria: 900 hours on each,
$900/6 = 150$ rolls and $900/12 = 75$ tonnes. Each consumes what it produces.

\textbf{4.} Cloth $75 + 150 = 225$; grain $150 + 75 = 225$.

\textbf{5.} Brenland in both goods (4 hours against 6, 2 against 12).

\textbf{6.} Kestria in cloth (\omterm{def:g10:ricardo-model:autarky-price}{autarky price} 0.5 against 2), Brenland in grain.

\textbf{7.} $0.5 < p < 2$.

\textbf{8.} Brenland only grain: 300 tonnes. Kestria only cloth: 300 rolls.

\textbf{9.} Brenland's output is worth 300 tonnes; half, 150 tonnes, is kept as grain and half buys $150/1 = 150$ rolls of
cloth. Kestria's output is worth $300 \times 1 = 300$ tonnes; it keeps 150 rolls and sells 150 for 150 tonnes of grain.
Each consumes 150 rolls and 150 tonnes.

\textbf{10.} Brenland exports 150 tonnes of grain and imports 150 rolls of cloth; Kestria the reverse. Cloth: 300 made,
$150 + 150 = 300$ consumed. Grain: 300 made, $150 + 150 = 300$ consumed. Both markets clear.

\textbf{11.} 300 rolls and 300 tonnes against 225 and 225: 75 more of each, a third more.

\textbf{12.} Brenland: 150 rolls and 150 tonnes against 75 and 150, a gain of 75 rolls of cloth. Kestria: 150 rolls and
150 tonnes against 150 and 75, a gain of 75 tonnes of grain.

\textbf{13.} Brenland gets a roll for one tonne of grain, 2 hours, instead of 4: it saves 2 hours per roll. Kestria gets a
tonne for one roll, 6 hours, instead of 12: it saves 6 hours per tonne.

\textbf{14.} A Brenland worker makes $1/2$ tonne an hour: $w_B = 0.5$ tonne. A Kestria worker makes $1/6$ roll an hour,
worth $1/6$ tonne at $p = 1$: $w_K = 1/6 \approx 0.167$ tonne. Ratio: $0.5 \times 6 = 3$.

\textbf{15.} A roll of cloth costs $6 \times \tfrac16 = 1$ tonne in Kestria and $4 \times 0.5 = 2$ tonnes in Brenland. A
tonne of grain costs $2 \times 0.5 = 1$ tonne in Brenland and $12 \times \tfrac16 = 2$ tonnes in Kestria. Each country
makes its export at half the other's cost.

\textbf{16.} Between $6/4 = 1.5$ (Brenland's lead in cloth) and $12/2 = 6$ (its lead in grain)
(\cref{prop:g10:ricardo-model:wage}).

\textbf{17.} Kestria's \omterm{def:g10:ricardo-model:autarky-price}{autarky price} of cloth becomes $6/8 = 0.75$ tonne, still below Brenland's 2: the pattern does not
change. Both specialise for $0.75 < p < 2$, and $p = 1$ is in this range.

\textbf{18.} Kestria still gets a tonne for one roll, 6 hours, but now could grow it in 8: it saves 2 hours per tonne
instead of 6. Trade saves the hours that the home method would have cost, and a better plough makes the home method
cheaper.

\textbf{19.} A plough that brings Kestria's hours per tonne down to 6: its \omterm{def:g10:ricardo-model:autarky-price}{autarky price} of cloth becomes $6/6 = 1 = p$,
and importing a tonne (6 hours of weaving) costs exactly what growing it does.

\textbf{20.} At $p = 1$, Brenland saves 2 hours of work on each roll of cloth it consumes, and Kestria 6 hours on each
tonne of grain; Brenland's wage is 3 times Kestria's, between the productivity ratios 1.5 and 6.
\end{solution}
